\section{Errors that build up and up}
\begin{quotation}\textit{
Up until perhaps the 1990s it was common to use floating point representations
that used 32 bits for each number. That corresponded to keeping around 7
decimal digits of precision. Despite that seeming generous it was fairly
easy to stumble across calculation schemes where the limited accuracy led
to errors building up and results ending up compromised. In those early
days floating point using 64-bit words (giving more like 16 decimal
digits) existed but their use was more expensive both in time and space.
When computer capacities had grown so that these limitations were not a
problem (at least for most users) it become more or less standard to use this
double precision. Effects of rounding felt a lot more remote with all that
extra precision in hand and so fewer people even thought about it.
However in the 2020s the pendulum swings back. Many of the most intense
calculations for both graphics and AI applications and performed not
using the main parts of a computer but exploiting the capabilities of
video cards or other special accelerators. These tend to favour low
precision calculations -- 32 bit floating point is at the high end of
normal usage and calculations are typically done using 16 bit or even
lower precisions. Furthermore the hardware is seriously fast so that
cascades of vast numbers of individual arithmetic operations will be
used. The effect is that maintaining accuracy comes back into focus.
Of course most ordinary users will not be coding up all these calculations
for themselves, and perhaps they do not mind too much if the results are
not especially accurate. But it may neverthless interest them to understand
some of the issues. In the coverage here almost all the examples will
use decimal fractions with three significant digits. This represents
a level of accuracy very close to that for 16 bit computer floating point
work, so any unreasonable-seeming results can transfer pretty directly to
that world.
}\end{quotation}



% add up n copies of 1/n
